% Acronyms
\usepackage{scalefnt,letltxmacro}
\LetLtxMacro{\oldtextsc}{\textsc}
\renewcommand{\textsc}[1]{\oldtextsc{\scalefont{1.10}#1}}
\usepackage[acronym,nowarn]{glossaries}
%\setacronymstyle{long-short-desc}

\glsdisablehyper
\makeglossaries
\usepackage{xspace}
\usepackage{float}
\usepackage{physics}
\usepackage[normalem]{ulem}

% Bullets
\usepackage{enumitem}
% Maths
\usepackage{amssymb}
\usepackage{mathtools}
\usepackage{amsfonts}
\usepackage{amsmath}
\usepackage{amsthm,thmtools}
\usepackage{booktabs}
% \usepackage[ruled,vlined]{algorithm2e}
\usepackage[]{microtype}

% Algorithms
\usepackage[linesnumbered,ruled,vlined]{algorithm2e}
\newcommand\mycommfont[1]{\footnotesize\ttfamily\textcolor{gray}{#1}}
\SetCommentSty{mycommfont}
\SetKwInput{KwInput}{Input}
\SetKwInput{KwOutput}{Output}

\usepackage{hyperref}
\usepackage{svg}

% cleverref
\usepackage[capitalize,nameinlink,]{cleveref} %noabbrev
\crefname{section}{\S}{\S\S}
\Crefname{section}{\S}{\S\S}
% \creflabelformat{equation}{#2\textup{#1}#3}
\providecommand\algorithmname{algorithm}

% color

%\definecolor{shadecolor}{gray}{0.9}
\newcommand{\red}[1]{\textcolor{BrickRed}{#1}}
\newcommand{\orange}[1]{\textcolor{BurntOrange}{#1}}
\newcommand{\green}[1]{\textcolor{OliveGreen}{#1}}
\newcommand{\blue}[1]{\textcolor{MidnightBlue}{#1}}
\newcommand{\sky}[1]{\textcolor{SkyBlue}{#1}}
\newcommand{\gray}[1]{\textcolor{black!60}{#1}}
%\definecolor{mylightgray}{gray}{0.94}

% counter for table
\makeatletter
\@ifundefined{c@rownum}{%
	\let\c@rownum\rownum
}{}
\@ifundefined{therownum}{%
	\def\therownum{\@arabic\rownum}%
}{}
\makeatother

% Other           
\DeclareRobustCommand{\parhead}[1]{\textbf{#1}~}

% For crossferencing across documents
\makeatletter
\newcommand*{\addFileDependency}[1]{% argument=file name and extension
	\typeout{(#1)}
	\@addtofilelist{#1}
	\IfFileExists{#1}{}{\typeout{No file #1.}}
}
\makeatother
\newcommand*{\myexternaldocument}[1]{%
	\externaldocument{#1}%
	\addFileDependency{#1.tex}%
	\addFileDependency{#1.aux}%
}
% % Figures
% \usepackage{nidanfloat}
\usepackage[font=footnotesize,labelfont=bf,tableposition=top]{caption}
\usepackage{tikz}
\usetikzlibrary{fit}
\usepackage{graphicx}
\usepackage[]{subcaption}   %% Already defined
% \usepackage{wrapfig}


\newcommand{\rulesep}{\unskip\ \hrule\ }

\usepackage{pgfplots}
\pgfplotsset{compat=1.6}
\usepgfplotslibrary{groupplots}
%% Use sans serif font for pictures
\tikzstyle{every picture}+=[font=\sffamily]
\tikzstyle{optimized} = [circle,fill=white,draw=black, dashed,inner sep=1pt, minimum size=20pt, font=\fontsize{10}{10}\selectfont, node distance=1]
\pgfkeys{/pgf/number format/.cd,1000 sep={}}
\pgfplotsset{
	tick label style = {font=\sffamily},
	every axis label/.append style={font=\sffamily},
	typeset ticklabels with strut,
}
\pgfplotsset{every axis/.append style={
			every x tick label/.append style={font=\fontsize{6pt}{6pt}\sffamily, yshift=.5ex,},
			every y tick label/.append style={font=\fontsize{6pt}{6pt}\sffamily, xshift=.5ex},
			every y label/.append style={xshift=10ex, font=\sffamily},
			every x label/.append style={yshift=3ex, font=\sffamily},
			every title/.append style={font=\sffamily}
		},
}
\pgfplotsset{
	xticklabel={$\mathsf{\pgfmathprintnumber{\tick}}$},
	yticklabel={$\mathsf{\pgfmathprintnumber{\tick}}$},
}
\pgfplotsset{every axis title/.append style={yshift=-1ex}}
\newlength\figureheight
\newlength\figurewidth
\usepgfplotslibrary{external}
\tikzexternalize[mode=list and make]
\tikzexternaldisable
\newcommand{\tikzfile}[1]{
	\tikzsetnextfilename{#1}%  
	\input{#1.tikz}
}
\newcommand*\circled[1]{\tikz[baseline=(char.base)]{
		\node[shape=circle,draw,inner sep=2pt] (char) {\tiny #1};}}




% \usepackage{physics}
% \usepackage{mathdots}
% \usepackage{yhmath}
% \usepackage{cancel}
% \usepackage{color}
% \usepackage{siunitx}
% \usepackage{array}
% \usepackage{multirow}
% \usepackage{amssymb}
% \usepackage{gensymb}
% \usepackage{tabularx}
% \usepackage{extarrows}
% \usepackage{booktabs}
% \usetikzlibrary{fadings}
% \usetikzlibrary{patterns}
% \usetikzlibrary{shadows.blur}
% \usetikzlibrary{shapes}


% ********************************************************************
% Notes and todos
% ********************************************************************
\usepackage[colorinlistoftodos,
	% disable,
	textsize=scriptsize,
	linecolor=red!30,
	bordercolor=red!30,
	backgroundcolor=red!10]{todonotes}
% \makeatletter
\renewcommand{\todo}[2][]{\tikzexternaldisable\@todo[#1]{#2}\tikzexternalenable}
% \makeatother
\newcommand{\note}[1]{\todo[linecolor=BurntOrange,backgroundcolor=BurntOrange!25,bordercolor=BurntOrange, inline, ]{\color{BurntOrange!80!black}{\textbf{\textit{#1}}}}}


\usepackage[most]{tcolorbox}
\newenvironment{highlight}%
{\begin{tcolorbox}[toprule=2mm,left=4pt,right=4pt,top=0pt,bottom=1pt,boxsep=2pt, before skip = 2ex, after skip =2ex]}{\end{tcolorbox}}
\tcbset{boxsep=4pt,left=2pt,right=2pt,top=-0pt,bottom=0pt}


\newcommand\noteMF[1]{{\color{red}{\footnotesize {\textbf{\textsc{MF}}: ``#1''}\xspace}}}
\newcommand\notePP[1]{{\color{blue}{\footnotesize {\textbf{\textsc{PP}}: {``#1''}}\xspace}}}
\newcommand\notePM[1]{{\color{orange}{\footnotesize {\textbf{\textsc{PM}}: ``#1''}\xspace}}}
\newcommand\noteGF[1]{{\color{violet}{\footnotesize {\textbf{\textsc{GF}}: ``#1''}\xspace}}}
\newcommand\noteEB[1]{{\color{green}{\footnotesize {\textbf{\textsc{EB}}: ``#1''}\xspace}}}
\newcommand\noteGC[1]{{\color{cyan}{\footnotesize {\textbf{\textsc{GC}}: ``#1''}\xspace}}}


\usepackage{url}

% \usepackage{xr}
% \externaldocument{supplementary}

%\usepackage{algorithm}
%\usepackage{algorithmic}
\usepackage{subcaption}

\usepackage{tikz}
\usepackage{graphicx}
%\usepackage{subfigure}
\usepackage{xcolor}
\usepackage{amsmath}
\usepackage{amssymb}
\usepackage{bm}
\usepackage{amsthm} % theorems and proofs
\usepackage{nicefrac}


\declaretheorem[name=Theorem]{theorem}
\declaretheorem[name=Lemma]{lemma}
\declaretheorem[name=Proposition]{proposition}
\declaretheorem[name=Corollary]{corollary}%numberwithin=section
\declaretheorem[name=Definition]{definition}
\declaretheorem[name=Condition]{condition}
\declaretheorem[name=Assumption]{assumption}
\declaretheorem[name=Example]{example}
\declaretheorem[name=Problem]{problem}


%\newtheorem{theorem}{Theorem}
%\newtheorem{proposition}{Proposition}
%\newtheorem{lemma}{Lemma}
%\newtheorem{sublemma}{Lemma}[lemma]
%\newtheorem{corollary}{Corollary}
%\newtheorem{definition}{Definition}

\usepackage{adjustbox}
\usepackage{multirow}
\usepackage{mathtools}
%\DeclareMathOperator{\Hessian}{H}
%\DeclareMathOperator{\Hessian}{\mathbf{H}}
\DeclareMathOperator{\Hessian}{Hess}

\usepackage{xspace}

%\hypersetup{draft}
\newcommand{\redquestionmark}{\textbf{\textcolor{red!60!black}{?}}}



